\chapter{Resultados en el conjunto sintético}
\label{chap:resultados}

El primer bloque de resultados establece si una CNN puede aprender la tarea en un entorno controlado. Este paso es necesario antes de evaluar el modelo sobre datos reales y antes de comparar con VLM/ICL. Si la CNN no logra aprender siquiera en un dominio sintético limpio, la comparación con VLM carecería de una línea base interpretable. El conjunto sintético contiene 100 pacientes simulados y 300 imágenes, con tres derivaciones por paciente. La etiqueta positiva corresponde a la presencia simultánea de RBBB, elevación ST e inversión de T en el simulador. Hay 20 pacientes positivos y 80 normales o incompletos.

\section{Curva por presupuesto}

La tabla \ref{tab:resultados_sinteticos_finales} muestra la media y la desviación estándar sobre tres semillas. La tendencia es clara: con \(K=2\) y \(K=4\) el modelo ya supera el azar en exactitud equilibrada, pero aún produce muchos falsos positivos. A partir de \(K=16\) aparece una mejora sustancial y con \(K=32\) se obtiene el mejor punto de la campaña.

\begin{table}[H]
  \centering
  \caption{Resultados CNN ResNet18 sobre el conjunto sintético QRS/ST.}
  \label{tab:resultados_sinteticos_finales}
  \small
  \begin{tabular}{rrrrrr}
    \toprule
    \(K\) & BA & F1 & Sens. & Esp. & ROC AUC \\
    \midrule
    2  & \(0,625 \pm 0,022\) & 0,401 & 0,650 & 0,600 & 0,520 \\
    4  & \(0,631 \pm 0,042\) & 0,407 & 0,617 & 0,646 & 0,566 \\
    8  & \(0,685 \pm 0,023\) & 0,468 & 0,683 & 0,688 & 0,629 \\
    16 & \(0,773 \pm 0,013\) & 0,569 & 0,800 & 0,746 & 0,755 \\
    32 & \(0,848 \pm 0,030\) & 0,673 & 0,883 & 0,812 & 0,857 \\
    \bottomrule
  \end{tabular}
\end{table}

\begin{figure}[H]
  \centering
  \safeincludegraphics[width=0.82\textwidth]{results/cnn_simulator_qrs_balanced_accuracy_by_k.png}
  \caption{Exactitud equilibrada de la CNN sobre el conjunto sintético según el presupuesto \(K\).}
  \label{fig:cnn_simulator_ba}
\end{figure}

El incremento entre \(K=8\) y \(K=16\) es el cambio más importante de la curva. El modelo pasa de una exactitud equilibrada media de 0,685 a 0,773, lo que indica que ocho ejemplos todavía no cubren suficientemente la variabilidad de combinaciones morfológicas. Con \(K=32\), la sensibilidad alcanza 0,883 y la especificidad 0,812, una combinación razonable para una tarea controlada de alerta morfológica.

Este resultado demuestra que la formulación visual es viable y que la CNN puede extraer información de las imágenes. Sin embargo, como se verá en el capítulo siguiente, este aprendizaje no se traslada al dominio real.

\section{Matriz de confusión}

La matriz agregada de \(K=32\) resume el mejor punto sintético. En las tres semillas se acumulan 53 verdaderos positivos, 195 verdaderos negativos, 45 falsos positivos y 7 falsos negativos. El error dominante sigue siendo marcar como positivos algunos pacientes incompletos, aunque el número de falsos negativos es bajo.

\begin{figure}[H]
  \centering
  \safeincludegraphics[width=0.58\textwidth]{results/cnn_simulator_qrs_k32_confusion_matrix.png}
  \caption{Matriz de confusión agregada del conjunto sintético para \(K=32\).}
  \label{fig:cnn_simulator_confusion}
\end{figure}

En un contexto de triaje, esta distribución de errores sería preferible a la situación inversa: el modelo tiende a avisar de más antes que a omitir positivos. Sin embargo, el conjunto sintético no representa pacientes reales y sus etiquetas proceden de reglas internas; por tanto, este resultado solo demuestra aprendizaje de la tarea simulada.

\section{Inspección cualitativa}

El panel Grad CAM de la figura \ref{fig:cnn_simulator_gradcam_resultados} muestra que el modelo atiende a regiones del complejo QRS, el punto J y la repolarización. La inspección es coherente con la definición de la etiqueta, pero no sustituye a una validación clínica. Grad CAM ayuda a detectar fallos obvios de atención; no certifica que la red razone como un especialista.

\begin{figure}[H]
  \centering
  \safeincludegraphics[width=0.92\textwidth]{results/cnn_simulator_qrs_gradcam_example.png}
  \caption{Ejemplo Grad CAM del modelo sintético.}
  \label{fig:cnn_simulator_gradcam_resultados}
\end{figure}

\section{Conclusión del dominio sintético}

El dominio sintético confirma tres hechos. Primero, el renderizado contiene información visual suficiente para entrenar una CNN con pocos ejemplos. Segundo, el presupuesto importa: por debajo de \(K=16\) el comportamiento es útil pero irregular; con \(K=32\) el modelo mejora simultáneamente sensibilidad, especificidad, F1 y área ROC. Tercero, este buen rendimiento no puede trasladarse de forma directa a clínica porque las imágenes sintéticas son más limpias y las etiquetas están definidas por el propio generador.

Este resultado no cierra la comparación con VLM/ICL. Establece que la CNN \emph{puede} aprender cuando las condiciones están controladas. La pregunta siguiente es si ese aprendizaje se mantiene en datos reales o si la CNN falla al enfrentarse a la variabilidad clínica. Si falla, la comparación con VLM/ICL se apoya en los propios datos.
